\documentclass[a4paper,twoside,11pt]{article}

\usepackage[french]{babel}
\usepackage[utf8]{inputenc}
\usepackage[T1]{fontenc}
\usepackage{lmodern}
\usepackage{amsmath,amssymb}
\usepackage{geometry}
\usepackage{minted}
\usepackage[most]{tcolorbox}
\usepackage{array}
\usepackage{hyperref}
\usepackage{fancyhdr}
\newcommand{\highlight}[1]{\textcolor{brown}{\textbf{#1}}}
\geometry{margin=2cm}
\hypersetup{
    colorlinks=true,
    linkcolor=brown,
    filecolor=magenta,      
    urlcolor=brown,
    pdftitle={Topologie 1},
    pdfpagemode=FullScreen,
}


% Styles pour les boîtes
\newtcolorbox{definitionbox}[1][]{
  colframe=cyan!20!cyan,
  colback=cyan!20,
  title={Définition},
  fonttitle=\bfseries,
  #1
}

\newtcolorbox{exemplebox}[1][]{
  title=\textbf{Exemple},
  colback=green!3,
  colframe=green!50!black,
  #1
}

\newtcolorbox{methodebox}[1][]{
  title=\textbf{Méthode / Analyse structurée},
  colback=orange!3,
  colframe=orange!70!black,
  #1
}

\newtcolorbox{attentionbox}[1][]{
  colback=red!20,
  colframe=red!50!red,
  title={Attention},
  fonttitle=\bfseries,
  #1
}


\newtcolorbox{profbox}[1][]{
  colframe=violet!20!violet,
  colback=violet!20,
  title={Exemple / Idée },
  fonttitle=\bfseries,
  #1
}

\newtcolorbox{redbox}[1][]{
  colback=red!20,
  colframe=red!50!red,
  title={Par coeur/ Attention},
  fonttitle=\bfseries,
  #1
}

\title{Analyse de la Topologie 4\\ Architecture Router-on-a-stick\large }
\author{florian mazé}
\begin{document}

\maketitle
\fancypagestyle{}{}
\pagestyle{fancy}
\fancyhead{} 
\fancyhead[RO,LE]{\textbf{Topologie 1}}
\fancyfoot{} 
\fancyfoot[LE,RO]{\thepage}
\fancyfoot[LO,CE]{florian mazé}
\fancyfoot[CO,RE]{}

\begin{center}
\includegraphics[scale=0.5]{"../images/topologie.png"}    
\end{center}

\tableofcontents

\section{Introduction}

\begin{profbox}[title=Compétences travaillées]
    \begin{itemize}
        \item Concevoir une solution d'infrastructure réseau (bloc 2)
        \item Installer tester et déployer une solution d'infrastructuré réseau (bloc 2)
        \item Sécuriser les équipements et les usages des utilisateurs (bloc 3)
    \end{itemize}
\end{profbox}


\begin{definitionbox}[breakable,title=Router-on-a-Stick (ROAS)]

Le \textit{Router-on-a-Stick} permet le routage inter-VLAN avec un seul lien physique entre un switch L2 et un routeur.\\
Le port du switch connecté au routeur est configuré en \texttt{trunk} (802.1Q) pour transporter plusieurs VLANs.\\
Côté routeur, on crée des \textit{sous-interfaces} sur l’interface physique (ex: \texttt{G0/0.10}, \texttt{G0/0.20}).\\
Chaque sous-interface est associée à un VLAN via \texttt{encapsulation dot1Q <vlan\_id>}.\\
On attribue une adresse IP à chaque sous-interface : c’est la passerelle par défaut des hôtes du VLAN.\\
Les machines de chaque VLAN utilisent cette IP comme \texttt{default gateway}.\\
Le switch conserve la commutation intra-VLAN, le routeur assure la communication entre VLANs.\\

\end{definitionbox}
\begin{redbox}
Avantage : simple et économique (pas besoin de switch L3 même si en l'occurrence j'en ai mis un).\\
Limite : le lien trunk et le routeur peuvent devenir un goulot d’étranglement.\\
\end{redbox}

\begin{profbox}[title=Améliorations à apporter]

  \begin{itemize}
      \item On sécurise souvent avec ACLs
      \item VLAN natif maîtrisé
      \item Trunk autorisant uniquement les VLANs nécessaires
  \end{itemize}

\end{profbox}



\section{Architecture de ségrégation réseau}

Dans le cadre de cette topologie, il s'agissait de représenter deux services de l'organisation:
\begin{itemize}
  \item Le service marketing 
  \item Le service des ressources humaines 
\end{itemize}
\bigskip 

\subsection{Création de VLANs}

Chaque service représente ici un segment logique du réseau physique:
\begin{itemize}
  \item Marketing - VLAN 22
  \item Ressouces Humaines - VLAN 20
\end{itemize}

\begin{definitionbox}[title=Définition d'un VLAN]
Un VLAN (\textit{Virtual LAN}) segmente logiquement un réseau sur une même infrastructure physique.\\
Chaque VLAN crée un domaine de broadcast distinct.\\
Deux hôtes dans le même VLAN communiquent directement via le switch (couche 2).\\
Deux hôtes dans des VLANs différents ne communiquent pas sans routage inter-VLAN.\\
Un switch peut porter plusieurs VLANs simultanément.\\
Un port \texttt{access} appartient à un seul VLAN (poste client, imprimante).\\
Un port \texttt{trunk} transporte plusieurs VLANs avec l’étiquetage 802.1Q.\\
Chaque VLAN est souvent associé à un sous-réseau IP différent.\\
La passerelle de chaque VLAN est fournie par un routeur ou un switch L3.\\
Objectifs : isolation, sécurité, performance et organisation du réseau.

\end{definitionbox}



\end{document}

